\section{Results}
\label{sec:results}

\subsection{Non-Compositional Compounds Are More Predictable}
\label{sec:results_farahmand}

\Tabref{tab:farahmand} presents the comparison between \comp and \noncomp noun compounds on the \farahmand dataset.
Across all metrics, \noncomp compounds show significantly higher within-phrase predictability.
The strongest effect appears in surprisal (Cohen's $d = -0.49$, $p < 0.0001$): \noncomp compounds have average surprisal of 7.63 bits compared to 9.46 bits for \comp compounds, a reduction of nearly two bits.
Next-token probability is also significantly higher for \noncomp compounds ($d = 0.21$, $p < 0.0001$).

The \logitlens signal is significant at all probed layers from 6 onward ($p < 0.005$), with effect sizes ranging from $d = 0.12$ (layer 12) to $d = 0.20$ (layers 6 and 42).
The \futurelens mid-layer probability from the first token shows a smaller but significant effect ($d = 0.12$, $p = 0.004$).

\begin{table}[t]
\centering
\caption{Comparison of prediction metrics between \comp (C, $n = 902$) and \noncomp (NC, $n = 140$) noun compounds on the \farahmand dataset. All $p$-values from two-sided Mann-Whitney $U$ tests. Best effect sizes are \textbf{bolded}.}
\label{tab:farahmand}
\vspace{4pt}
\resizebox{\textwidth}{!}{%
\begin{tabular}{@{}lcccc@{}}
\toprule
Metric & C (mean $\pm$ sd) & NC (mean $\pm$ sd) & $p$-value & Cohen's $d$ \\
\midrule
NTP avg probability & $0.028 \pm 0.090$ & $0.049 \pm 0.110$ & $< 0.0001$ & $0.21$ \\
NTP avg surprisal & $9.46 \pm 3.86$ & $7.63 \pm 3.64$ & $< 0.0001$ & $\mathbf{-0.49}$ \\
\midrule
\logitlens L6 & $0.011 \pm 0.052$ & $0.027 \pm 0.102$ & $< 0.0001$ & $0.20$ \\
\logitlens L12 & $0.015 \pm 0.067$ & $0.028 \pm 0.100$ & $0.002$ & $0.12$ \\
\logitlens L24 & $0.026 \pm 0.099$ & $0.045 \pm 0.130$ & $0.001$ & $0.16$ \\
\logitlens L42 & $0.048 \pm 0.144$ & $0.077 \pm 0.178$ & $< 0.0001$ & $0.18$ \\
\logitlens L47 & $0.028 \pm 0.094$ & $0.049 \pm 0.138$ & $< 0.0001$ & $0.18$ \\
\midrule
\futurelens mid prob & $0.024 \pm 0.094$ & $0.036 \pm 0.114$ & $0.004$ & $0.12$ \\
\futurelens L47 $k{=}1$ & $0.028 \pm 0.094$ & $0.049 \pm 0.138$ & $< 0.0001$ & $0.18$ \\
\bottomrule
\end{tabular}%
}
\end{table}

\subsection{Idioms Show Dramatically Higher Predictability}
\label{sec:results_idioms}

\Tabref{tab:idioms} compares \idiomem idioms against \comp control bigrams.
The effects are very large by conventional standards.
Next-token probability within idioms averages 0.277 compared to 0.021 for control bigrams (Cohen's $d = 1.85$, $p < 0.0001$).
The \logitlens at the final layer shows a similarly large effect ($d = 1.76$).
Once inside an idiom, the model is highly confident about subsequent tokens---a clear signature of memorization as a holistic unit.

\begin{table}[t]
\centering
\caption{Comparison of prediction metrics between \comp control bigrams ($n = 55$) and \idiomem idioms ($n = 852$). Effect sizes above 1.0 are \textbf{bolded}.}
\label{tab:idioms}
\vspace{4pt}
\resizebox{\textwidth}{!}{%
\begin{tabular}{@{}lcccc@{}}
\toprule
Metric & Control (mean $\pm$ sd) & Idioms (mean $\pm$ sd) & $p$-value & Cohen's $d$ \\
\midrule
NTP avg probability & $0.021 \pm 0.032$ & $0.277 \pm 0.193$ & $< 0.0001$ & $\mathbf{1.85}$ \\
NTP avg surprisal & $7.34 \pm 2.54$ & $5.56 \pm 2.07$ & $< 0.0001$ & $-0.77$ \\
\midrule
\logitlens L42 & $0.050 \pm 0.108$ & $0.232 \pm 0.182$ & $< 0.0001$ & $\mathbf{1.22}$ \\
\logitlens L47 & $0.024 \pm 0.037$ & $0.268 \pm 0.193$ & $< 0.0001$ & $\mathbf{1.76}$ \\
\bottomrule
\end{tabular}%
}
\end{table}

\subsection{From-First-Token Prediction Behaves Differently for Long Phrases}
\label{sec:results_futurelens}

A notable finding is that the \futurelens from-first-token signal behaves oppositely for idioms compared to noun compounds.
Idioms show \emph{lower} \futurelens probability from the first token at middle layers than control bigrams ($p_{\text{FL}}^{\text{mid}} = 0.002$ vs.\ $0.019$, $p < 0.0001$).
This occurs because idioms are longer (3--5 tokens), making it harder for a single hidden state to encode the entire continuation.
The within-phrase prediction signals (NTP, \logitlens)---which measure prediction at each position given full context up to that point---remain highly effective regardless of phrase length.

This result highlights an important methodological point: the \emph{within-phrase} prediction signal is more robust for detecting memorization than the \emph{from-first-token} signal, especially for longer phrases.

\subsection{Correlation with Human Compositionality Scores}
\label{sec:results_correlation}

\Tabref{tab:correlation} presents Spearman correlations between prediction metrics and graded non-compositionality scores on the \farahmand dataset.
All correlations are statistically significant ($p < 0.0001$) and positive, indicating that higher prediction accuracy corresponds to lower compositionality.
The strongest correlations are for NTP surprisal ($r = -0.24$) and NTP average probability ($r = 0.23$), followed by \logitlens at the final layer ($r = 0.22$).

\begin{table}[t]
\centering
\caption{Spearman correlation between prediction metrics and non-compositionality score (fraction of judges, 0--1) on the \farahmand dataset ($n = 1{,}042$). All $p < 0.0001$.}
\label{tab:correlation}
\vspace{4pt}
\begin{tabular}{@{}lc@{}}
\toprule
Metric & Spearman $r$ \\
\midrule
NTP avg probability & $\mathbf{0.233}$ \\
NTP avg surprisal & $-0.235$ \\
\logitlens L47 & $0.224$ \\
\logitlens L42 & $0.209$ \\
\futurelens L47 $k{=}1$ & $0.224$ \\
\futurelens mid prob & $0.149$ \\
\bottomrule
\end{tabular}
\end{table}

\subsection{Binary Classification Performance}
\label{sec:results_classification}

\Tabref{tab:classification} reports ROC AUC for binary classification of \noncomp vs.\ \comp compounds.
The best single metric is NTP average probability (AUC = 0.643), slightly outperforming a 5-fold cross-validated logistic regression model over all 22 features (AUC = $0.619 \pm 0.011$).
This suggests the features are correlated and provide limited complementary information when combined.

When pooling all \noncomp phrases (Farahmand NC + idioms) against all \comp phrases (Farahmand C + control bigrams), NTP average probability achieves AUC = 0.899.
However, this comparison is confounded by phrase type and length differences and should be interpreted with caution.

\begin{table}[t]
\centering
\caption{ROC AUC for binary non-compositionality classification on the \farahmand dataset.}
\label{tab:classification}
\vspace{4pt}
\begin{tabular}{@{}lc@{}}
\toprule
Method & AUC \\
\midrule
NTP avg probability & $\mathbf{0.643}$ \\
NTP avg surprisal & $0.640$ \\
\logitlens L47 & $0.633$ \\
\logitlens L42 & $0.624$ \\
\futurelens L47 $k{=}1$ & $0.632$ \\
\futurelens mid layer & $0.575$ \\
\midrule
Logistic regression (22 features, 5-fold CV) & $0.619 \pm 0.011$ \\
\bottomrule
\end{tabular}
\end{table}

\subsection{Layer-Wise Profile of the Non-Compositional Advantage}
\label{sec:results_layers}

The \logitlens layer profile reveals distinct processing patterns across phrase types.
Control bigrams maintain low prediction probability across all layers.
Farahmand compounds show slightly higher prediction, with \noncomp compounds consistently above \comp ones at every layer from 6 onward (all $p < 0.005$).
Idioms show a distinctive trajectory: very low prediction at early layers (0--18), followed by a sharp increase at layers 24+, reaching 0.27 at layer 47.

This profile is consistent with the literature on transformer processing: early layers operate on individual tokens, while deeper layers construct holistic phrase-level representations.
The sharp late-layer rise for idioms suggests that memorized phrase retrieval is a deep-layer phenomenon.
